% META: {"id": "TSPE-LOGARITHME-QCM", "chapitre": "TSPE-LOGARITHME", "type_objet": "qcm", "genere_depuis": "chapitres/TSPE-LOGARITHME/qcm/TSPE-LOGARITHME-QCM.json", "status": "generated"}
% Fichier genere par scripts/build_qcm_tex.py — ne pas editer a la main.

\section*{\textcolor{chapcolor}{\MakeUppercase{QCM d'auto-evaluation — Fonction logarithme népérien}}}

\begin{center}
\textit{Pour chaque question, une seule reponse est exacte.}
\end{center}

\begin{enumerate}

\item[] \textbf{Capacite C1}

\item \textbf{[Q1]} $\ln(2x)$ (pour $x>0$) est egal a :
  \begin{enumerate}[label=\Alph*.]
    \item $2\ln x$.
    \item $\ln 2 + \ln x$.
    \item $\ln 2 \times \ln x$.
    \item $(\ln x)^2$.
  \end{enumerate}

\item \textbf{[Q2]} La derivee de $\ln(5x+1)$ est :
  \begin{enumerate}[label=\Alph*.]
    \item $\dfrac{1}{5x+1}$.
    \item $\dfrac{5}{5x+1}$.
    \item $\dfrac{5x+1}{5}$.
    \item $5\ln(5x+1)$.
  \end{enumerate}

\item \textbf{[Q3]} $\displaystyle\lim_{x \to +\infty} \ln x$ vaut :
  \begin{enumerate}[label=\Alph*.]
    \item $0$.
    \item $-\infty$.
    \item $+\infty$.
    \item $1$.
  \end{enumerate}

\item \textbf{[Q4]} $\displaystyle\lim_{x \to 0^+} x\ln x$ vaut :
  \begin{enumerate}[label=\Alph*.]
    \item $-\infty$.
    \item $0$.
    \item $1$.
    \item Cette limite n'existe pas.
  \end{enumerate}

\item \textbf{[Q5]} Pour resoudre $\mathrm{e}^{2x} = 7$, on applique aux deux membres :
  \begin{enumerate}[label=\Alph*.]
    \item la fonction carre.
    \item la fonction $\ln$.
    \item la fonction exponentielle a nouveau.
    \item l'inverse.
  \end{enumerate}

\end{enumerate}

\ifnxVersionProfesseur
\clearpage
\section*{Cle de correction — reservee au professeur}

\begin{center}
\begin{tabular}{lll}
\hline
Question & Capacite & Reponse exacte \\
\hline
Q1 & C1 & \textbf{B} \\
Q2 & C1 & \textbf{B} \\
Q3 & C1 & \textbf{C} \\
Q4 & C1 & \textbf{B} \\
Q5 & C1 & \textbf{B} \\
\hline
\end{tabular}
\end{center}

\subsection*{Diagnostic des reponses erronees}

\begin{itemize}
  \item \textbf{Q1}
  \begin{itemize}
    \item \textbf{A} — Le coefficient $2$ est a l'interieur du logarithme : la propriete donne $\ln(2x)=\ln 2+\ln x$, pas $2\ln x$. \emph{Renvoi : C1, logarithme d'un produit.}
    \item \textbf{C} — Le logarithme transforme un produit en somme, pas en produit de logarithmes. \emph{Renvoi : C1, logarithme d'un produit.}
    \item \textbf{D} — Aucune propriete ne transforme $\ln(2x)$ en $(\ln x)^2$ ; il faut separer les deux facteurs. \emph{Renvoi : C1, logarithme d'un produit.}
  \end{itemize}
  \item \textbf{Q2}
  \begin{itemize}
    \item \textbf{A} — L'eleve a utilise $(\ln u)'=1/u$ sans multiplier par $u'=5$. \emph{Renvoi : C1, derivee de ln(u).}
    \item \textbf{C} — La formule est $u'/u$, et non $u/u'$ : le numerateur doit etre la derivee $5$. \emph{Renvoi : C1, derivee de ln(u).}
    \item \textbf{D} — Deriver une composee logarithmique ne consiste pas a multiplier le logarithme par la derivee interieure. \emph{Renvoi : C1, derivee de ln(u).}
  \end{itemize}
  \item \textbf{Q3}
  \begin{itemize}
    \item \textbf{A} — $\ln x$ ne tend pas vers $0$ lorsque $x$ tend vers $+\infty$ ; sa croissance est lente mais non bornee. \emph{Renvoi : C1, limites du logarithme.}
    \item \textbf{B} — $\ln x$ est croissante sur $]0,+\infty[$ ; elle ne peut pas tendre vers $-\infty$ a droite. \emph{Renvoi : C1, limites du logarithme.}
    \item \textbf{D} — $1$ est la valeur $\ln(\mathrm e)$, pas la limite de $\ln x$ en $+\infty$. \emph{Renvoi : C1, limites du logarithme.}
  \end{itemize}
  \item \textbf{Q4}
  \begin{itemize}
    \item \textbf{A} — Bien que $\ln x$ tende vers $-\infty$, le facteur $x$ tend vers $0$ assez vite pour que le produit tende vers $0$. \emph{Renvoi : C1, croissance comparee.}
    \item \textbf{C} — La limite remarquable $x\ln x$ en $0^+$ vaut $0$, et non $1$. \emph{Renvoi : C1, croissance comparee.}
    \item \textbf{D} — La limite existe : en posant $x=1/t$, on obtient $-\ln(t)/t$, qui tend vers $0$. \emph{Renvoi : C1, croissance comparee.}
  \end{itemize}
  \item \textbf{Q5}
  \begin{itemize}
    \item \textbf{A} — La fonction carre ne transforme pas directement l'equation exponentielle en une equation lineaire en $x$. \emph{Renvoi : C1, equations exponentielles.}
    \item \textbf{C} — Appliquer encore l'exponentielle complique l'equation ; $\ln$ est la fonction reciproque adaptee. \emph{Renvoi : C1, equations exponentielles.}
    \item \textbf{D} — Prendre les inverses donne $\mathrm e^{-2x}=1/7$ sans isoler directement $x$ ; appliquer $\ln$ donne $2x=\ln 7$. \emph{Renvoi : C1, equations exponentielles.}
  \end{itemize}
\end{itemize}
\fi
